%===============================================================================================
% PER-VARIANT VISUAL IDENTITY.
%===============================================================================================
% Section ORDER alone was never enough. Six documents with identical fonts, name treatment,
% section rules, bullet glyph and spacing read as one template with the blocks shuffled — which
% was exactly the complaint. Each block below overrides the universal defaults so that every
% variant is a different document at a glance, before a word of it is read.
%
% THE AXES, AND WHY THESE ONES. Every axis fires for EVERY profile. Nothing here is contingent on
% optional data, so no variant can collapse onto its neighbour because a field happened to be
% empty — the way D used to become byte-identical to C whenever a project had no tech stack.
%
%   variant   section header        name block            bullet      skills          density
%   -------   -------------------   -------------------   ---------   -------------   --------
%   A         rule beneath          centred small caps    en dash     label per line  normal
%   B         rules above + below   centred + underline   disc        label table     airy
%   C         small caps, no rule   centred bold roman    asterisk    flowing prose   tight
%   D         vertical bar, left    left-aligned block    open square two columns     normal
%   E         shaded bar            centred, wide rule    open circle hanging indent  airy
%   F         short accent rule     name over full rule   none        label table     tight
%
% Experience entry shape is a seventh axis, chosen in the content builder (VARIANT_STYLE) because
% it decides which command to call: inline for A, C and F, stacked-company for B and E,
% stacked-role for D.
%
% CONSTRAINTS HONOURED THROUGHOUT:
%   * No package beyond what _preamble.tex already loads.
%   * No colour. Variant E's bar is achromatic grey via the `gray` model, which is the one
%     shaded treatment the design brief sanctioned; it is a filled box behind live text, not an
%     image, so extraction is unaffected.
%   * Single column throughout. Variant D's two-column block is multicol around the skills list
%     only — short label/value pairs, which extract in reading order. The document body is never
%     multi-column, because that is what actually breaks an ATS.
%   * Base font size is never reduced below 10pt. Density is spacing, not shrinking.
%
% CONTAINS NO CANDIDATE DATA.
%===============================================================================================

%%STYLE a%%
% ---- A: the reference look. Bold section label over a full-width rule, centred small-caps name,
%         en-dash bullets, one skills line per category, normal density, one-line entries. Its
%         distinguishing mark against the rest is the pipe-separated skills list.
% The below-space has to absorb the \vspace{-5pt} the header's after-code applies once the rule is
% drawn; at 4pt the net was negative and the first line of every section sat ON the rule.
\titlespacing*{\section}{0pt}{9pt}{10pt}
\setlist[itemize]{topsep=9pt, partopsep=0pt, parsep=0pt, itemsep=1pt}

%%STYLE b%%
% ---- B: rules ABOVE and below a small-caps label — the section headers read as banded strips
%         rather than underlined words. Airy: the most generous spacing of the six. Skills sit in
%         a two-column label table, and experience entries stack company over an italic role.
% \titleline puts the rule on its OWN full-width line. A bare \titlerule in the before-code fills
% the line itself and shoves the label to the right margin, which is how the first attempt at this
% ended up with right-aligned section names.
\titleformat{\section}[block]{\raggedright\normalsize\scshape\bfseries}{}{0em}{\titleline{\titlerule}\vspace{3pt}}[\vspace{2pt}\titleline{\titlerule}]
\titlespacing*{\section}{0pt}{11pt}{8pt}
\setlist[itemize]{topsep=13pt, partopsep=0pt, parsep=0pt, itemsep=4pt}

\renewcommand{\resumebullet}{\textbullet}

\renewcommand{\resumeHeadline}[1]{\textbf{\small #1}\\ \vspace{3pt}}
\renewcommand{\resumeHeader}[3]{
\begin{center}
{\LARGE \bfseries #1} \\ \vspace{4pt}
#2#3
\end{center}
\vspace{-6pt}
}

% PLAIN tabular, not tabularx. tabularx grabs its whole body as a macro ARGUMENT so it can
% measure the X column, which means it cannot find its own \end when the begin and end live in
% separate commands — it aborts while still scanning for the end of its own body. Two fixed
% p-columns need no measuring pass and split across macros cleanly.
\renewcommand{\resumeSkillsStart}{\begingroup\small\begin{tabular}{@{}p{0.24\textwidth}p{0.72\textwidth}@{}}}
\renewcommand{\resumeSkillEntry}[2]{\textbf{#1} & #2 \\[2pt]}
\renewcommand{\resumeSkillsEnd}{\end{tabular}\endgroup\vspace{-2pt}}

%%STYLE c%%
% ---- C: no rules anywhere. Section labels are small caps carried by whitespace alone, which is
%         the quietest and most typographic of the six. Tight: the densest spacing, so a long
%         profile fits a page. Skills run as one flowing paragraph with bold labels.
\titleformat{\section}{\raggedright\large\scshape\bfseries}{}{0em}{}[\vspace{-2pt}]
\titlespacing*{\section}{0pt}{9pt}{2pt}
\setlist[itemize]{topsep=8pt, partopsep=0pt, parsep=0pt, itemsep=0pt}

\renewcommand{\resumebullet}{\textasteriskcentered}

\renewcommand{\resumeHeader}[3]{
\begin{center}
{\Huge \bfseries #1} \\ \vspace{2pt}
#2#3
\end{center}
\vspace{-10pt}
}

% One flowing paragraph, categories separated by a spaced glyph. The separator LEADS each entry
% and is suppressed for the first, rather than trailing each entry — a trailing separator leaves a
% dangling glyph after the last category, and \unskip removes the space but not the glyph.
% \newif is plain TeX; this is preamble scope, so it is legal here and needs no package.
\newif\ifresumeskillfirst
\renewcommand{\resumeSkillsStart}{\begingroup\small\noindent\resumeskillfirsttrue}
\renewcommand{\resumeSkillEntry}[2]{\ifresumeskillfirst\resumeskillfirstfalse\else\hspace{6pt}\resumebullet\hspace{6pt}\fi\textbf{#1:} #2}
\renewcommand{\resumeSkillsEnd}{\par\endgroup\vspace{-2pt}}

%%STYLE d%%
% ---- D: the only left-aligned masthead, and the only one whose section labels are marked by a
%         vertical bar instead of a horizontal rule. Skills run in two columns. Every one of those
%         fires on every profile — which is the point, because D's old sole differentiator (the
%         project tech stack on its own line) vanished for any profile that recorded no tech
%         stack, leaving D byte-identical to C. It keeps that tech-stack treatment as a bonus.
\titleformat{\section}{\raggedright\large\bfseries}{}{0em}{\rule[-0.55ex]{2.5pt}{1.6ex}\hspace{7pt}}[\vspace{-4pt}]
\titlespacing*{\section}{0pt}{10pt}{4pt}
\setlist[itemize]{topsep=10pt, partopsep=0pt, parsep=0pt, itemsep=1pt}

\renewcommand{\resumebullet}{\scalebox{0.62}{$\Box$}}
\renewcommand{\resumeEntrySecondary}[1]{\small#1}

% Left-aligned block: name, then headline and contact beneath it, all flush left, closed by a
% rule. Nothing is centred, which is the difference visible from across the room.
\renewcommand{\resumeHeadline}[1]{\textbf{\small #1}\\ \vspace{2pt}}
\renewcommand{\resumeHeader}[3]{
\noindent{\LARGE \bfseries #1} \\ \vspace{3pt}
\noindent #2#3
\vspace{3pt}
\noindent\rule{\linewidth}{0.4pt}
\vspace{-4pt}
}

\renewcommand{\resumeSkillsStart}{\begingroup\small\begin{multicols}{2}}
\renewcommand{\resumeSkillEntry}[2]{\noindent\textbf{#1:} #2\par\vspace{2pt}}
\renewcommand{\resumeSkillsEnd}{\end{multicols}\endgroup\vspace{-2pt}}

%%STYLE e%%
% ---- E: section labels sit in a shaded bar running the full text width — the heaviest header
%         treatment of the six and unmistakable at a glance. Airy spacing, open-circle bullets,
%         and a hanging indent on the skills lines so a long category wraps under its own text
%         rather than back to the margin.
\newcommand{\resumeShadedTitle}[1]{\colorbox[gray]{0.90}{\makebox[\dimexpr\textwidth-2\fboxsep\relax][l]{\strut\normalsize\bfseries #1}}}
\titleformat{\section}{\raggedright}{}{0em}{\resumeShadedTitle}[\vspace{-2pt}]
\titlespacing*{\section}{0pt}{12pt}{6pt}
\setlist[itemize]{topsep=12pt, partopsep=0pt, parsep=0pt, itemsep=3pt}

\renewcommand{\resumebullet}{\scalebox{0.75}{$\circ$}}

\renewcommand{\resumeHeadline}[1]{\textbf{\small #1}\\ \vspace{3pt}}
\renewcommand{\resumeHeader}[3]{
\begin{center}
{\LARGE \scshape #1} \\ \vspace{3pt}
#2#3
\end{center}
\vspace{-6pt}
}

\renewcommand{\resumeSkillsStart}{\begingroup\small}
\renewcommand{\resumeSkillEntry}[2]{\noindent\hangindent=1.4em\hangafter=1\textbf{#1:} #2\par\vspace{3pt}}
\renewcommand{\resumeSkillsEnd}{\par\endgroup\vspace{-3pt}}

%%STYLE f%%
% ---- F: a short accent rule under each section label rather than a full-width one, and the name
%         set over a rule spanning the page. Tight spacing, and the only variant with NO bullet
%         glyph at all — bullets are hanging-indented lines, which reads as a report rather than a
%         list. Each \item still starts its own line in the PDF, so extraction is unaffected.
\titleformat{\section}{\raggedright\large\bfseries}{}{0em}{}[\vspace{-7pt}\rule{0.16\textwidth}{1.3pt}\vspace{-3pt}]
\titlespacing*{\section}{0pt}{10pt}{3pt}
\setlist[itemize]{topsep=8pt, partopsep=0pt, parsep=0pt, itemsep=0pt}

\renewcommand{\resumebullet}{}
\renewcommand{\resumeItemListStart}{\begin{itemize}[leftmargin=1.1em, label={}, itemindent=0pt]}

\renewcommand{\resumeHeadline}[1]{\textbf{\small #1}\\ \vspace{2pt}}
\renewcommand{\resumeHeader}[3]{
\begin{center}
{\LARGE \bfseries #1}
\end{center}
\vspace{-6pt}
\noindent\rule{\linewidth}{0.4pt}
% center adds its own \topsep above, so this pulls the block back under the rule rather than
% leaving a band of white between the name and the contact line.
\vspace{-16pt}
\begin{center}
#2#3
\end{center}
\vspace{-10pt}
}

\renewcommand{\resumeSkillsStart}{\begingroup\small\begin{tabular}{@{}p{0.22\textwidth}p{0.74\textwidth}@{}}}
\renewcommand{\resumeSkillEntry}[2]{\textbf{#1} & #2 \\}
\renewcommand{\resumeSkillsEnd}{\end{tabular}\endgroup\vspace{-2pt}}
